\section{Formula\c{c}\~ao Determin\'istica Equivalente}
\label{sec:deq}

Na Formula\c{c}\~ao Determin\'istica Equivalente (DEQ), o problema estoc\'astico multiestágio \'e representado explicitamente sobre a \'arvore de cen\'arios. Cada n\'o $n \in \mathcal{N}$ corresponde a um estado poss\'ivel do sistema em determinado est\'agio, incorporando o hist\'orico de realiza\c{c}\~oes das incertezas at\'e aquele ponto. As decis\~oes de contrata\c{c}\~ao s\~ao, portanto, indexadas por n\'o, e a estrutura de n\~ao antecipatividade \'e imposta explicitamente por meio de restri\c{c}\~oes de igualdade entre n\'os irm\~aos.

\subsection{N\~ao Antecipatividade}
\label{subsec:nao_antecipatividade}

Para garantir que a comercializadora n\~ao utilize informa\c{c}\~ao futura em suas decis\~oes, imp\~oe-se que os volumes negociados sejam id\^enticos para todos os n\'os que compartilham o mesmo predecessor. Para cada contrato $i \in \mathcal{I}$:
%
\begin{equation}
  q^B_{i,m} = q^B_{i,n} \quad \text{e} \quad q^S_{i,m} = q^S_{i,n},
  \qquad \forall\, m, n \in \mathcal{N} \mid a(m) = a(n).
  \label{eq:nao_antecipatividade_deq}
\end{equation}

Os volumes respeitam os limites de lote ($0 \leq q^B_{i,n} \leq \bar{q}^B_i$) e s\~ao nulos quando o n\'o n\~ao pertence ao m\^es de in\'icio do contrato ($q^B_{i,n} = 0$, $\forall\, i \notin \mathcal{I}_n$). O racioc\'inio \'e id\^entico para $q^S_{i,n}$.

\subsection{Din\^amica do Portf\'olio}
\label{subsec:dinamica_portfolio}

Contratos com dura\c{c}\~ao superior a um m\^es geram efeitos financeiros e f\'isicos ao longo de v\'arios est\'agios. O \'indice $k$ representa o n\'umero de meses remanescentes de vig\^encia de um contrato. As equa\c{c}\~oes de recurs\~ao s\~ao escritas em termos de $D_{\max}$ para generalidade; neste trabalho, $D_{\max} = 6$ meses \'e o limite pr\'atico imposto na confer\^encia de contratos. A evolu\c{c}\~ao da posi\c{c}\~ao f\'isica l\'iquida $v_{s,k,n}$ \'e dada por:
%
\begin{align}
  v_{s,k,n} &= v_{s,k+1,a(n)} + \sum_{\substack{i \in \mathcal{I}_n \\ d_i > k,\; s_i = s}}
  \!\bigl(q^B_{i,n} - q^S_{i,n}\bigr),
  &\forall\, s \in \mathcal{S},\; n \in \mathcal{N} \setminus \{n_0\},\; k < D_{\max},
  \label{eq:volume_fisico} \\[4pt]
  v_{s,D_{\max},n} &= \sum_{\substack{i \in \mathcal{I}_n \\ d_i > D_{\max},\; s_i = s}}
  \!\bigl(q^B_{i,n} - q^S_{i,n}\bigr),
  &\forall\, s \in \mathcal{S},\; n \in \mathcal{N} \setminus \{n_0\}.
  \label{eq:volume_fisico_max}
\end{align}

De forma an\'aloga, a posi\c{c}\~ao financeira l\'iquida $c_{s,k,n}$ evolui segundo:
%
\begin{align}
  c_{s,k,n} &= c_{s,k+1,a(n)} + \sum_{\substack{i \in \mathcal{I}_n \\ d_i > k,\; s_i = s}}
  \!\bigl(q^S_{i,n} K^S_i - q^B_{i,n} K^B_i\bigr),
  &\forall\, s \in \mathcal{S},\; n \in \mathcal{N} \setminus \{n_0\},\; k < D_{\max},
  \label{eq:posicao_financeira} \\[4pt]
  c_{s,D_{\max},n} &= \sum_{\substack{i \in \mathcal{I}_n \\ d_i > D_{\max},\; s_i = s}}
  \!\bigl(q^S_{i,n} K^S_i - q^B_{i,n} K^B_i\bigr),
  &\forall\, s \in \mathcal{S},\; n \in \mathcal{N} \setminus \{n_0\}.
  \label{eq:posicao_financeira_max}
\end{align}

\subsection{Exposi\c{c}\~ao ao Mercado de Curto Prazo}
\label{subsec:exposicao}

Ap\'os a realiza\c{c}\~ao das incertezas no n\'o $n$, calcula-se a exposi\c{c}\~ao da comercializadora ao mercado de curto prazo. A vari\'avel $E_{s,n}$ consolida a gera\c{c}\~ao realizada, a carteira legada e os contratos vigentes:
%
\begin{equation}
  E_{s,n} = G_{s,n} + V^{0B}_{s,n} - V^{0S}_{s,n}
  + \sum_{\substack{i \in \mathcal{I}_n \\ s_i = s}} \!\bigl(q^B_{i,n} - q^S_{i,n}\bigr)
  + v_{s,1,a(n)},
  \qquad \forall\, s \in \mathcal{S},\; n \in \mathcal{N} \setminus \{n_0\}.
  \label{eq:exposicao}
\end{equation}

O termo $v_{s,1,a(n)}$ representa as posi\c{c}\~oes f\'isicas herdadas do n\'o predecessor $a(n)$ com exatamente $k = 1$ m\^es remanescente de vig\^encia --- ou seja, contratos que encerram no per\'iodo corrente. Posi\c{c}\~oes com $k \geq 2$ ainda t\^em meses futuros de vig\^encia e ser\~ao incorporadas \`a exposi\c{c}\~ao dos est\'agios seguintes por meio da recurs\~ao~\eqref{eq:volume_fisico}.

\subsection{Equa\c{c}\~ao de Transi\c{c}\~ao de Caixa}
\label{subsec:transicao_caixa}

O saldo de caixa $x_n$ agrega o saldo herdado do predecessor, o resultado da carteira legada, o fluxo dos novos contratos, o fluxo dos contratos j\'a vigentes herdados, e a liquida\c{c}\~ao da exposi\c{c}\~ao ao PLD. O termo $R^{\text{leg}}_n$ representa exclusivamente o fluxo financeiro a pre\c{c}o fixo da carteira legada ---  isto \'e, $\sum_{s} (K^{0S}_{s,n} V^{0S}_{s,n} - K^{0B}_{s,n} V^{0B}_{s,n}) h_n$ --- sem incluir a liquida\c{c}\~ao ao PLD. A exposi\c{c}\~ao ao mercado de curto prazo dos volumes legados \'e capturada integralmente pela vari\'avel $E_{s,n}$ definida em~\eqref{eq:exposicao}, que j\'a cont\'em $V^{0B}_{s,n} - V^{0S}_{s,n}$; desse modo, o resultado econ\^omico total de um contrato legado --- diferencial entre o pre\c{c}o contratado e o PLD realizado --- \'e corretamente contabilizado sem dupla contagem:
%
\begin{equation}
  x_n = x_{a(n)} + R^{\text{leg}}_n
  + \left[\sum_{i \in \mathcal{I}_n}\!\bigl(q^S_{i,n} K^S_i - q^B_{i,n} K^B_i\bigr)
  + \sum_{s \in \mathcal{S}} c_{s,1,a(n)}\right] h_n
  + \sum_{s \in \mathcal{S}} E_{s,n}\, P_{s,n}\, h_n,
  \label{eq:transicao_caixa}
\end{equation}
%
para todo $n \in \mathcal{N} \setminus \{n_0\}$.

As dimens\~oes s\~ao consistentes: $q^S_{i,n} K^S_i$ tem unidade MWm~$\times$~R\$/MWh = R\$/h; multiplicado por $h_n$~h, resulta em R\$. De forma an\'aloga, $c_{s,1,a(n)}$ em R\$/h vezes $h_n$ resulta em R\$; e $E_{s,n} P_{s,n}$ em MWm~$\times$~R\$/MWh = R\$/h vezes $h_n$ resulta em R\$.

A restri\c{c}\~ao de cr\'edito \'e imposta de forma el\'astica por meio de uma vari\'avel de folga $\delta_n \geq 0$, preservando a factibilidade do problema:
%
\begin{equation}
  x_n + \delta_n \geq L^{\text{cred}}, \qquad \forall\, n \in \mathcal{N} \setminus \{n_0\}.
  \label{eq:restricao_credito}
\end{equation}
%
Neste trabalho, $L^{\text{cred}} = -10^8$~R\$ e $M = 10^9$~R\$. O valor de $M$ \'e escolhido como \emph{big-M}: suficientemente grande para que qualquer viola\c{c}\~ao do limite seja penalizada al\'em do resultado financeiro m\'aximo ating\'ivel pelo portf\'olio, sem introduzir instabilidade num\'erica.

\subsection{Fun\c{c}\~ao Objetivo}
\label{subsec:objetivo_deq}

O modelo maximiza o resultado financeiro esperado ao final do horizonte, descontando penalidades associadas \`as viola\c{c}\~oes do limite de cr\'edito:
%
\begin{equation}
  \max \quad \sum_{n \in \mathcal{L}} \pi_n\, x_n
  \;-\; M \sum_{n \in \mathcal{N} \setminus \{n_0\}} \pi_n\, \delta_n.
  \label{eq:objetivo_deq}
\end{equation}
%
A penalidade incide sobre todos os n\'os $n \in \mathcal{N} \setminus \{n_0\}$, e n\~ao apenas sobre as folhas, pois a restri\c{c}\~ao de cr\'edito~\eqref{eq:restricao_credito} deve ser satisfeita em todos os est\'agios do horizonte --- refletindo um requisito de solv\^encia intra-horizonte.

A formula\c{c}\~ao resultante \'e um problema de programa\c{c}\~ao linear (PL) cujo porte cresce com o n\'umero de n\'os da \'arvore de cen\'arios. Para $R$ ramos por est\'agio e $T$ est\'agios, o n\'umero de n\'os \'e da ordem de $O(R^T)$, o que torna a abordagem computacionalmente invi\'avel para inst\^ancias de maior dimens\~ao, como demonstrado nos experimentos do Cap\'itulo~\ref{cap:experimentos}.
